\chapter{自同态的约化}\label{ch:b2:reduction}

要理解一个自同态，只需找出它仅仅
拉伸的方向。本章建立这套工具 --- 特征值、
特征多项式与极小多项式、核分解
引理 --- 以及它的收获：对角化与三角化
的判别准则、凯莱--哈密顿定理、邓福德分解，以及
计算幂与指数的方法，它们将为
\cref{ch:b2:diffeq}所用。全文中 $E$ 是一个
有限维 $K$-向量空间（$K = \R$ 或 $\C$），且 $u \in
\mathcal{L}(E)$, $n = \dim E$.

\section{特征值与特征向量}

\begin{definition}\label{def:b2:reduction:eigen}
$\lambda \in K$ 称为\emph{特征值}\index{特征值}，即对 $u$ 而言，当
$u(x) = \lambda x$ 对某个 $x \neq 0$ 成立时（
即\emph{特征向量}）；\emph{特征子空间}为 $E_\lambda(u) =
\ker(u - \lambda\,\mathrm{id})$。特征值的全体称为
\emph{谱}\index{谱} $\operatorname{Sp}(u)$。子空间 $F$
称为\emph{稳定的}，当 $u(F) \subseteq F$ 时；特征子空间都是稳定的，并且
稳定子空间允许定义诱导自同态 $u|_F$。
\end{definition}

\begin{theorem}[特征子空间的独立性]\label{thm:b2:reduction:independence}
属于两两不同特征值的特征向量构成一个
线性无关的向量族；等价地，特征子空间的和
$E_{\lambda_1} + \dots + E_{\lambda_r}$（$\lambda_i$ 两两不同）是
直和。特别地，$u$ 至多有 $n$ 个特征值。
\end{theorem}

\begin{proof}
对 $r$ 作归纳。设 $x_1 + \dots + x_r = 0$，其中 $x_i \in
E_{\lambda_i}$，已知结论对 $r - 1$ 成立。两边作用 $u$，
并减去 $\lambda_r$ 倍的原关系式：
\[
\sum_{i=1}^{r-1} (\lambda_i - \lambda_r)\, x_i = 0 ,
\]
于是由归纳假设，逐项有 $(\lambda_i - \lambda_r)x_i = 0$，即 $x_i =
0$（$i < r$ 时），于是 $x_r = 0$。
维数为 $n$ 的空间中，非零子空间的直和至多含 $n$ 个加项。
\end{proof}

\begin{omfigure}
\begin{tikzpicture}[scale=0.9, >={Stealth}]
  \draw[gray!40] (-2.6,-2.2) grid (2.6,2.2);
  \draw[->] (-2.6,0) -- (2.7,0) node[below, font=\small] {$x$};
  \draw[->] (0,-2.2) -- (0,2.3) node[left, font=\small] {$y$};
  \draw[omDef, very thick, ->] (0,0) -- (0.7,0.7)
    node[above, font=\small] {$v_1$};
  \draw[omDef, thick, ->, dashed] (0,0) -- (2.1,2.1)
    node[above right, font=\small] {$Av_1 = 3v_1$};
  \draw[omThm, very thick, ->] (0,0) -- (0.9,-0.9)
    node[below, font=\small] {$v_2$};
  \draw[omThm, thick, ->, dashed] (0,0) -- (0.9,-0.9);
  \draw[omProp, very thick, ->] (0,0) -- (1,0)
    node[below right, font=\small] {$e_1$};
  \draw[omProp, thick, ->, dashed] (0,0) -- (2,1)
    node[right, font=\small] {$Ae_1$};
\end{tikzpicture}

{\small 矩阵 $A = \begin{psmallmatrix}2 & 1\\ 1 &
2\end{psmallmatrix}$ 作用在平面上：一般向量
$e_1$ 被推离了它所在的直线，而特征方向 $v_1 =
(1,1)$ 与 $v_2 = (1,-1)$ 仅被拉伸 --- 分别伸缩 $3$ 倍与
$1$ 倍（故 $Av_2 = v_2$：虚线的像与 $v_2$ 重合）。
对角化即换到基 $(v_1, v_2)$，此时
$A$ 化为 $\operatorname{diag}(3, 1)$。}
\end{omfigure}

\begin{definition}[特征多项式]\label{def:b2:reduction:charpoly}
$\chi_u(X) = \det(X\,\mathrm{id} - u)$\index{特征多项式
} --- 在任一基下计算为 $\det(XI_n - A)$，
是 $n$ 次首一多项式，在相似变换下不变
（见\cref{thm:b2:linalg:detrules}）。其在 $K$ 中的根恰好是
特征值（$\lambda$ 是特征值 $\iff u - \lambda\,\mathrm{id}$
不是单射 $\iff \chi_u(\lambda) = 0$），并且
\[
\chi_u(X) = X^n - (\operatorname{tr} u)\, X^{n-1} + \dots +
(-1)^n \det u .
\]
特征值的\emph{代数重数}\index{代数重数}
$m_\lambda$ 是它作为 $\chi_u$ 的根的
重数；\emph{几何
重数}\index{几何重数}为 $\dim E_\lambda$，且 $1 \leq \dim E_\lambda \leq
m_\lambda$.
\end{definition}

\begin{proof}[所述事实的证明]
系数部分：用置换公式展开 $\det(XI - A)$；
恒等置换贡献 $\prod_i (X - a_{ii})
= X^n - (\sum a_{ii})X^{n-1} + \dots$，其余每个置换
至多固定 $n - 2$ 个对角位置，贡献的次数 $\leq
n - 2$：最高两项的系数如上所述；令 $X = 0$ 得到
常数项 $\det(-A) = (-1)^n\det A$。

几何重数 $\leq$ 代数重数：设 $d = \dim E_\lambda$，将
$E_\lambda$ 的一组基扩充为 $E$ 的基，此时 $u$ 的矩阵是
分块上三角，左上块为 $\lambda I_d$，故 $\chi_u(X) =
(X - \lambda)^d\, \chi_{\text{（下方块）}}(X)$：
$\lambda$ 的重数至少为 $d$。
\end{proof}

\begin{example}[相同 $\chi$，几何不同]\label{ex:b2:reduction:samechi}
下列两个矩阵
\[
\begin{pmatrix}2 & 0\\ 0 & 2\end{pmatrix}
\qquad\text{和}\qquad
\begin{pmatrix}2 & 1\\ 0 & 2\end{pmatrix}
\]
它们有相同的特征多项式 $(X - 2)^2$、迹、
行列式与谱 --- 但二者不相似：第一个矩阵的
$E_2$ 的维数为 $2$（几何重数为 $2$），第二个矩阵
的维数为 $1$。特征多项式只能看到
代数重数；特征子空间的维数是
更精细的不变量，而极小多项式来裁决（$X - 2$
与 $(X - 2)^2$）。对所有可对角化性的讨论，其要点是：
$\chi$ 负责圈定候选者，而
由核来投票。
\end{example}

\begin{definition}[可对角化，可三角化]\label{def:b2:reduction:diag}
称 $u$ 为\emph{可对角化}\index{可对角化}，当 $E$ 存在
由特征向量组成的基（矩阵语言：与某个对角矩阵相似）；
\emph{可三角化}\index{可三角化}，当它在某组
基下的矩阵为上三角矩阵。
\end{definition}

\begin{theorem}[可对角化的判别准则]\label{thm:b2:reduction:diagcrit}
下列命题等价：
\begin{enumerate}
  \item $u$ 可对角化；
  \item $E = \bigoplus_{\lambda \in \operatorname{Sp} u} E_\lambda$;
  \item $\chi_u$ 在 $K$ 上分裂，且 $\dim E_\lambda = m_\lambda$
        对每个特征值成立；
  \item （充分但不必要）$\chi_u$ 有 $n$ 个
        互不相同的根，且均在 $K$ 中。
\end{enumerate}
\end{theorem}

\begin{proof}
（1 $\iff$ 2）：特征向量构成的基可分成
各 $E_\lambda$ 的基；反之，把直和的各个
直和项的基拼接起来即得 $E$ 的一组基
（见\cref{thm:b2:reduction:independence}，该和为直和；
维数相等则和即为全部）。

（2 $\iff$ 3）：在对角基下，$\chi_u = \prod (X -
\lambda)^{\dim E_\lambda}$ 分裂，且重数相符。
反之，设 $\chi_u$ 分裂，且 $\dim E_\lambda =
m_\lambda$ 对所有特征值成立，则特征子空间的
直和（由\cref{thm:b2:reduction:independence}知其为直和）的维数为
\[
\sum_{\lambda}\dim E_\lambda = \sum_{\lambda} m_\lambda =
\deg\chi_u = n ,
\]
中间的等号成立，因为分裂多项式的次数恰是
其根的重数之和；于是该和即为整个 $E$。这里要注意
两个假设各自的作用：分裂用来凑足次数，
重数相等用来凑足维数。

（4 $\Rightarrow$ 1）：$n$ 个互异特征值给出 $n$ 个线性无关的
特征向量（见\cref{thm:b2:reduction:independence}）：构成一组基。
\end{proof}

\begin{method}[判定可对角化]\label{met:b2:reduction:decide}
实际中按以下顺序检验 --- 每一步都可能
直接完成判定。（1）是否自然出现根为单根且完全分裂的
零化多项式（$u^2 = \mathrm{id}$，$u^2 = u$，$u^k =
\mathrm{id}$）？若是，则可对角化，无需计算
（见下文\cref{cor:b2:reduction:minpolycrit}）。（2）计算
$\chi_u$；若它有 $n$ 个互异的根，且它们均属于 $K$，
则可对角化（见\cref{thm:b2:reduction:diagcrit}~(4)）。（3）否则，仅对每个
重根 $\lambda$，比较 $\dim\ker(u -
\lambda\,\mathrm{id})$ 与重数 $m_\lambda$：任何
亏缺即排除可对角化；处处相等则证明之。
切勿计算单根的特征子空间（其维数必定
为 $1$），也切勿仅为判定而三角化。
\end{method}

\begin{example}[对角化的应用]\label{ex:b2:reduction:diagwork}
$A = I + J = \begin{psmallmatrix}2 & 1 & 1\\ 1 & 2 & 1\\ 1 & 1 &
2\end{psmallmatrix}$，其中 $J$ 为全一矩阵：由
$\operatorname{Sp}(J) = \{3, 0\}$
（见\cref{ex:b2:linalg:onesmatrix}），得 $\operatorname{Sp}(A) = \{4,
1\}$，特征子空间为 $\R(1,1,1)$ 与平面 $\{x + y + z =
0\}$：维数分别为 $1 + 2 = 3$，可对角化
（见\cref{thm:b2:reduction:diagcrit}~(2)）。求幂无需任何
过渡矩阵：取 $\Pi = J/3$，它是投影到
$\R(1,1,1)$,
\[
A = 4\,\Pi + 1\cdot(I - \Pi)
\quad\Longrightarrow\quad
A^k = 4^k\,\Pi + (I - \Pi)
= \frac{4^k - 1}{3}\,J + I .
\]
（验算 $k = 1$：$\frac{4-1}3 J + I = A$。）要点是：
当特征子空间清晰可见时，谱\emph{投影算子}
计算幂的速度比 $PDP^{-1}$ 所能达到的还快 --- 而且
公式还直观展示了动态：$A^k$ 以 $4^k$ 的速率
沿 $(1,1,1)$ 方向增长，而在正交平面上保持不变。
\end{example}

\begin{theorem}[三角化]\label{thm:b2:reduction:trigonalization}
$u$ 在 $K$ 上可三角化，当且仅当 $\chi_u$ 在
$K$ 上分裂。特别地，每个 $\C$ 上向量空间的自同态都是
可三角化。\index{三角化}
\end{theorem}

\begin{proof}
（$\Rightarrow$）三角矩阵的特征多项式
为 $\prod(X - t_{ii})$，因而分裂。

（$\Leftarrow$）对 $n$ 作归纳。因 $\chi_u$ 分裂，它有一个
根 $\lambda$，取特征向量 $e_1$。在以
$e_1$ 为首的基下，矩阵为 $\begin{pmatrix} \lambda & \ast\\ 0 &
B\end{pmatrix}$，且 $\chi_u = (X - \lambda)\chi_B$：$\chi_B$ 也
分裂。对 $(n-1) \times
(n-1)$ 矩阵 $B$ 应用归纳假设，存在可逆的 $Q$，
使 $Q^{-1}BQ$ 为上三角矩阵；以 $\begin{pmatrix}1 & 0\\
0 & Q\end{pmatrix}$ 对整个矩阵作共轭，即将其三角化。
\end{proof}

\begin{example}[手工三角化]\label{ex:b2:reduction:trigwork}
$B = \begin{pmatrix}3 & -1\\ 1 & 1\end{pmatrix}$：$\chi_B = X^2 -
4X + 4 = (X - 2)^2$，且 $\ker(B - 2I) =
\ker\begin{psmallmatrix}1 & -1\\ 1 & -1\end{psmallmatrix}$ 为
由 $e_1' = (1, 1)$ 张成的直线：只有一个特征值，
其特征子空间为一维 --- 不可对角化，但
可三角化（见\cref{thm:b2:reduction:trigonalization}）。
以 $e_2' = (1, 0)$ 补全为基，并计算：
\[
u(e_1') = (2, 2) = 2e_1',
\qquad
u(e_2') = (3, 1) = 1\cdot e_1' + 2\, e_2' ,
\]
于是在基 $(e_1', e_2')$ 下，矩阵为 $T =
\begin{psmallmatrix}2 & 1\\ 0 & 2\end{psmallmatrix}$。要点是：
$T$ 的对角元是被迫的（两个元素都
必须是二重特征值 $2$）；只有右上角的元素
依赖于 $e_2'$ 的选取，而重新缩放 $e_2'$ 即可使它
取任一非零值 --- 那个顽固的“$1$”是
邓福德将要分离出的幂零部分的影子。
\end{example}

\section{自同态的多项式}

\begin{definition}\label{def:b2:reduction:polyu}
设 $P = \sum a_k X^k \in K[X]$，记 $P(u) = \sum a_k u^k \in
\mathcal{L}(E)$。映射 $P \mapsto P(u)$ 是
$K[X] \to \mathcal{L}(E)$ 的代数同态（\cref{def:b2:structures:algebra}）；
其核 $\{P : P(u) = 0\}$ 是 $K[X]$ 的理想，且非零（
因为族 $(\mathrm{id}, u, \dots, u^{n^2})$ 在
$n^2$ 维的 $\mathcal{L}(E)$ 中线性相关），故它由唯一的
首一多项式 $\mu_u$ 生成。称之为\emph{极小
多项式}\index{极小多项式}
（\cref{thm:b2:structures:principal}）。
\end{definition}

\begin{proposition}\label{prop:b2:reduction:minpoly}
\begin{enumerate}
  \item $P(u) = 0 \iff \mu_u \mid P$；$u$ 的特征值是
        每个零化多项式的根，且 $\mu_u$ 的根
        \emph{恰为}特征值。
  \item 若 $F$ 是不变子空间，则 $\mu_{u|_F} \mid \mu_u$。
\end{enumerate}
\end{proposition}

\begin{proof}
(1) 整除性即生成元的定义。若 $u(x) =
\lambda x$，$x \neq 0$，则 $0 = P(u)(x) = P(\lambda)x$，于是
$P(\lambda) = 0$：特征值是零化多项式的根，
特别是 $\mu_u$ 的根。反之，若 $\lambda$ 为根，则 $\mu_u =
(X - \lambda)Q$，且 $Q(u) \neq 0$（$\mu_u$ 的次数极小）。取
$y$ 使 $Q(u)(y) \neq 0$，则 $(u - \lambda)(Q(u)(y)) = \mu_u(u)(y)
= 0$，故 $Q(u)(y)$ 即为特征向量。

(2) $\mu_u(u|_F) = \mu_u(u)|_F = 0$，并对 $u|_F$ 应用 (1)。
\end{proof}

\begin{example}[手工求极小多项式]\label{ex:b2:reduction:minpolyhand}
极小多项式通过逐次检验
次数来确定。对全一矩阵 $J \in \mathcal{M}_3(\R)$：
$J \neq \lambda I$（次数 $1$ 已排除），且 $J^2 = 3J$，于是
\[
\mu_J = X^2 - 3X = X(X - 3) :
\]
次数为 $2$，分裂，根为单根 --- $J$ 可对角化，谱为
$\{0, 3\}$（见下文\cref{cor:b2:reduction:minpolycrit}），
印证了\cref{ex:b2:linalg:onesmatrix}，无需
任何行列式计算。对
\cref{ex:b2:reduction:projectorswork}中的交换矩阵 $A$：$A \neq \pm I$ 且 $A^2
= I$ 得 $\mu_A = X^2 - 1$。两种情形的规律相同：
即从结构猜出一个低次恒等式
（秩一导出 $J^2 = (\operatorname{tr}J)\,J$；对合导出
$A^2 = I$），再验证没有真因子
能零化它。极小多项式通常是\emph{找出}的，而非
由 $\chi$ 计算出的。
\end{example}

\begin{theorem}[核分解引理]\label{thm:b2:reduction:kernels}
若 $P = P_1 P_2 \cdots P_r$，且各 $P_i$ 两两互素，
则\index{核分解引理}
\[
\ker P(u) = \ker P_1(u) \oplus \dots \oplus \ker P_r(u),
\]
且到各直和项的投影都是 $u$ 的多项式。
\end{theorem}

\begin{proof}
只需对 $r = 2$ 证明，再用归纳法。由 $K[X]$ 中的贝祖
恒等式（见\cref{thm:b2:structures:principal}）：$U P_1 + V P_2 = 1$，于是
对每个 $x$，有
\[
x = \underbrace{U(u)P_1(u)(x)}_{=:\,x_2}
+ \underbrace{V(u)P_2(u)(x)}_{=:\,x_1}.
\]
若 $x \in \ker P(u)$：$P_2(u)(x_2) = U(u)\,P(u)(x) = 0$（
$u$ 的多项式彼此可交换），故 $x_2 \in \ker P_2(u)$，对称地 $x_1
\in \ker P_1(u)$。两直和项之和充满 $\ker P(u)$；二者都
位于 $\ker P(u)$ 内（$P_i \mid P$）。直和性：若 $x \in \ker P_1(u)
\cap \ker P_2(u)$，则 $x = U(u)P_1(u)x + V(u)P_2(u)x = 0$。
$x_1, x_2$ 的公式给出投影为 $V(u)P_2(u)$ 与
$U(u)P_1(u)$。
\end{proof}

\begin{example}[显式投影的核引理]\label{ex:b2:reduction:projectorswork}
设 $A = \begin{psmallmatrix}0 & 1 & 0\\ 1 & 0 & 0\\ 0 & 0 &
1\end{psmallmatrix}$（交换前两个坐标）。于是 $A^2
= I$：多项式 $X^2 - 1 = (X - 1)(X + 1)$ 零化 $A$，
其因子互素，且贝祖恒等式是显式的：
\[
\frac{1}{2}(X + 1) - \frac12(X - 1) = 1 .
\]
依照\cref{thm:b2:reduction:kernels}的证明，
到 $\ker(A - I)$ 与 $\ker(A + I)$ 上的投影是
$A$ 的多项式：
\[
\pi_+ = \frac{A + I}{2} = \frac12\begin{pmatrix}
1 & 1 & 0\\ 1 & 1 & 0\\ 0 & 0 & 2\end{pmatrix},
\qquad
\pi_- = \frac{I - A}{2} = \frac12\begin{pmatrix}
1 & -1 & 0\\ -1 & 1 & 0\\ 0 & 0 & 0\end{pmatrix}.
\]
验证：$\pi_+ + \pi_- = I$，$\pi_+\pi_- = 0$，$\pi_\pm^2 =
\pi_\pm$，且像分别是平面 $\{x = y\}$（对称向量，特征值
$1$）与直线 $\R(1, -1, 0)$
（反对称，特征值 $-1$）。核引理并非存在性陈述：
贝祖系数\emph{即为}
投影公式。
\end{example}

\begin{example}[投影子同样可用来计算指数]\label{ex:b2:reduction:expprojectors}
同一交换矩阵，更进一步。由于 $A = \pi_+ -
\pi_-$，这两个投影子在代数意义上
正交（$\pi_+\pi_- = 0$）。每个幂次均满足 $A^k = \pi_+ +
(-1)^k\pi_-$，指数级数按投影子重新组合：
\[
\eu^{tA} = \sum_k \frac{t^k}{k!}\bigl(\pi_+ +
(-1)^k\pi_-\bigr)
= \eu^{t}\,\pi_+ + \eu^{-t}\,\pi_- =
\begin{pmatrix}
\cosh t & \sinh t & 0\\
\sinh t & \cosh t & 0\\
0 & 0 & \eu^{t}
\end{pmatrix}.
\]
（验证：$t = 0$ 时为单位矩阵；$0$ 处导数为 $A$。）
特征分解把一个矩阵级数化为两个
数项级数 --- 正是\cref{ch:b2:diffeq}将在每个可对角化
方程组上运用的精确机制，也是双曲函数
支配对称耦合的原因。
\end{example}

\begin{corollary}[借助极小多项式判定可对角化]\label{cor:b2:reduction:minpolycrit}
$u$ 可对角化 $\iff$ $\mu_u$ 在 $K$ 上分裂，
且根均为\emph{单根}，$\iff$ $u$ 的某个零化多项式
分裂且根为单根。
\end{corollary}

\begin{proof}
若 $P(u) = 0$，其中 $P = \prod_{i}(X - \lambda_i)$（
$\lambda_i$ 两两互异），由引理得 $E = \ker P(u) = \bigoplus_i \ker(u -
\lambda_i)$：即特征子空间的直和，故 $u$ 可对角化
（见\cref{thm:b2:reduction:diagcrit}）。反之，可对角化的
$u$ 被 $\prod_{\lambda \in \operatorname{Sp}u}(X -
\lambda)$ 零化（它零化每个特征子空间），该多项式分裂且
单根；$\mu_u$ 整除它，且二者有相同的根
（见\cref{prop:b2:reduction:minpoly}），故 $\mu_u$ 恰为该乘积。
\end{proof}

\begin{example}\label{ex:b2:reduction:projections}
投影满足 $p^2 = p$：被 $X(X-1)$ 零化，分裂且根为
单根 --- 可对角化，谱 $\subseteq \{0, 1\}$，且 $E =
\ker p \oplus \ker(p - \mathrm{id})$：本科一年级的几何分析，
于一行内重新证明。对称变换（$s^2 = \mathrm{id}$，
零化子为 $X^2 - 1$）：当 $\operatorname{char} K
\neq 2$ 时可对角化，谱 $\subseteq \{\pm 1\}$。一个满足 $u^3 =
u^2$ 且 $u^2 \neq u$：被 $X^2(X - 1)$ 零化，\emph{未必}
可对角化 --- 判别准则能检测出这一点（二重根
$0$ 须加以检验：可对角化当且仅当还有 $\ker u^2 = \ker
u$）。
\end{example}

\begin{example}[域决定结论：$\R^3$ 中的旋转]\label{ex:b2:reduction:rotation3d}
设 $R$ 为绕 $z$ 轴的四分之一转：
\[
R = \begin{pmatrix}
0 & -1 & 0\\
1 & 0 & 0\\
0 & 0 & 1
\end{pmatrix},
\qquad
\chi_R = (X - 1)(X^2 + 1).
\]
在 $\R$ 上：唯一特征值为 $1$，特征子空间为
$\R e_3$ --- 即不动向量构成的一条直线，再无进一步的
化简：$R$ 既不可对角化，也不可三角化，
在 $\mathcal{M}_3(\R)$ 中（$\chi_R$ 不分裂）。在 $\C$ 上有三个
互异特征值 $1, \iu, -\iu$，故 $R$ 可对角化，
特征向量为 $e_3$ 与 $e_1 \mp \iu e_2$。几何图景在代数中已有
反映：平面内的旋转没有
实不变方向，而复特征值 $\pm\iu$
（模为 $1$）储存了角度（$\pm\frac\pi2$），
这是实矩阵只能通过混合坐标才能表达的。
\end{example}

\begin{example}[极小多项式与特征多项式]\label{ex:b2:reduction:minvschar}
设 $D = \operatorname{diag}(2, 2, 3)$：$\chi_D = (X - 2)^2(X -
3)$，但 $\mu_D = (X - 2)(X - 3)$，因为 $(D - 2I)(D - 3I) = 0$
（在标准基上验证），而单独的任一因子都不零化
$D$。对于移位块 $N = \begin{psmallmatrix}0 & 1\\ 0 &
0\end{psmallmatrix} \oplus (3)$，即 $N' =
\begin{psmallmatrix}0 & 1 & 0\\ 0 & 0 & 0\\ 0 & 0 &
3\end{psmallmatrix}$：$\chi_{N'} = X^2(X - 3)$ \emph{且}
$\mu_{N'} = X^2(X - 3)$ --- 二重根确有必要，
因为 $N'$ 在 $\ker$ 一侧不可对角化（$N'e_2 =
e_1 \neq 0$）。经验法则：$\mu$ 与 $\chi$ 有相同的根
（见\cref{prop:b2:reduction:minpoly}）；$\mu$ 中的重数
刻画最大幂零块的大小，
$\chi$ 中的重数则是特征子空间的总维数。
\end{example}

\begin{theorem}[凯莱--哈密顿]\label{thm:b2:reduction:cayleyhamilton}
$\chi_u(u) = 0$；因此 $\mu_u \mid \chi_u$，且 $\deg \mu_u
\leq n$。\index{凯莱--哈密顿定理}
\end{theorem}

\begin{proof}
取 $x \neq 0$，令 $d$ 为极大的整数，使得 $(x, u(x), \dots,
u^{d-1}(x))$ 线性无关；记
\[
u^d(x) = -a_0 x - a_1 u(x) - \dots - a_{d-1}u^{d-1}(x),
\]
并令 $P_x = X^d + a_{d-1}X^{d-1} + \dots + a_0$，于是 $P_x(u)(x) =
0$。将该线性无关组扩充为 $E$ 的一组基，在此基下 $u$ 具有
分块形式 $\begin{pmatrix} C & \ast\\ 0 & D\end{pmatrix}$，其中 $C$
为 $P_x$ 的伴侣矩阵，其特征多项式为
$P_x$（沿第一列展开 $\det(XI - C)$，对
$d$ 用归纳法）。因此 $\chi_u = P_x \cdot \chi_D$，且
\[
\chi_u(u)(x) = \chi_D(u)\bigl(P_x(u)(x)\bigr) = 0 .
\]
对每个 $x$ 该论证都成立，故 $\chi_u(u) = 0$。
\end{proof}

\begin{example}[凯莱--哈密顿定理的应用]\label{ex:b2:reduction:chwork}
设 $A = \begin{pmatrix} 1 & 2\\ 3 & 4\end{pmatrix}$：$\chi_A = X^2 -
5X - 2$，故 $A^2 = 5A + 2I$。$A$ 的任意幂都可化为
$I$ 与 $A$ 的线性组合：
\[
A^4 = (5A + 2I)^2 = 25A^2 + 20A + 4I = 145A + 54I =
\begin{pmatrix} 199 & 290\\ 435 & 634\end{pmatrix},
\]
逆矩阵随之而来：由 $A(A - 5I) = 2I$ 可得
\[
A^{-1} = \tfrac12(A - 5I)
= \begin{pmatrix} -2 & 1\\ 3/2 & -1/2\end{pmatrix}.
\]
小结：凯莱--哈密顿定理把整个代数
$K[A]$ 压缩到 $\operatorname{Vect}(I, A, \dots, A^{n-1})$ ---
$\dim K[A] = \deg\mu_A \leq n$，无论需要的幂次有多大。
\end{example}

\begin{remark}[常见误区]\label{rem:b2:reduction:pitfalls}
(i) 特征值不可相加：$\operatorname{Sp}(A + B)$ 并不等于
$\operatorname{Sp}A + \operatorname{Sp}B$，且可对角化矩阵的和
未必可对角化 ---
$\begin{psmallmatrix}1 & 1\\ 0 & 0\end{psmallmatrix} +
\begin{psmallmatrix}0 & 0\\ 0 & 1\end{psmallmatrix} =
\begin{psmallmatrix}1 & 1\\ 0 & 1\end{psmallmatrix}$ 是两个可对角化矩阵之和
（各自有互异特征值），而
它本身不可对角化；只有\emph{可交换}族才表现良好
（见\cref{exo:b2:reduction:9}）。(ii) “$\chi_u$ 分裂”是关于
\emph{域}的假设：平面旋转的特征多项式为 $\chi =
X^2 - 2\cos\theta\,X + 1$，在 $\C$ 上分裂，而不在 $\R$ 上 ---
在 $\mathcal{M}_2(\C)$ 中可对角化，但
在 $\mathcal{M}_2(\R)$ 中不可三角化。(iii) 不等式的方向是几何重数 $\leq$
代数重数，绝不能颠倒；仅检验 $\dim E_\lambda \geq
1$ 并不能推出可对角化。(iv) $\mu_u$ 并不等于
$\chi_u$：等号成立当且仅当每个特征值恰有
一条块链（例如伴侣矩阵，以及本章的
周末问题）；以 $\chi$ 代替所需的 $\mu$，
会使每一次幂运算都变得繁重。(v) 邓福德的 $d$ 与 $\nu$ 是
$u$ 的多项式 ——一个具有适当性质的分解 $u = d' + \nu'$，
但 $d'\nu' \neq \nu'd'$，则\emph{并非}邓福德分解，
且绝不唯一。
\end{remark}

\begin{remark}[本章的应用]
约化是全书其余部分的主力：矩阵的幂与
矩阵指数驱动了\cref{ch:b2:diffeq}中的线性微分方程组；
\cref{ch:b2:quadratic}中的谱定理是
正交化的对角化；生成函数
（\cref{ch:b2:genfun}）以解析方法重新推导本章
周末问题的递推渐近性态。在本科三年级分册中，
同样的纲领在无穷维中展开：
紧自伴算子的谱理论，其中特征值数列取代了
有限谱，以及正矩阵的佩龙–弗罗贝尼乌斯理论，
它解释了\emph{为何}计数问题的主特征值
是正的且单的。
\end{remark}

\section{幂零自同态与邓福德分解}

\begin{proposition}[幂零自同态]\label{prop:b2:reduction:nilpotent}
对 $u$，若 $\chi_u$ 分裂，则下列各条件等价：$u^n =
0$；$u^k = 0$ 对某个 $k$ 成立；$\operatorname{Sp}(u) = \{0\}$；$\chi_u
= X^n$；$u$ 可三角化且对角线全为零。幂零
自同态的极小多项式为 $\mu_u = X^{\text{(幂零指数)}}$，且指数
$\leq n$。\index{幂零自同态}
\end{proposition}

\begin{proof}
$u^k = 0$ 使每个特征值都是 $X^k$ 的根，谱为 $\{0\}$
（当 $\chi$ 分裂时谱非空 ——在 $\C$ 上总成立）。于是 $\chi_u =
X^n$（根全为零），由凯莱–哈密顿定理得 $u^n = 0$；
三角化（\cref{thm:b2:reduction:trigonalization}）使
对角线全为零（对角线恰为特征值）。
反之，设 $A$ 为严格上三角矩阵：$a_{ij} = 0$
对 $j \leq i$ 成立。下面用归纳法证明
\[
(A^k)_{ij} = 0 \qquad \text{当 } j \leq i + k - 1,
\]
即每次乘幂都把零区域向上推移一条对角线。
当 $k = 1$ 时即为假设。对于归纳步，
\[
(A^{k+1})_{ij} = \sum_{\ell} (A^k)_{i\ell}\,a_{\ell j} ,
\]
各项皆为零：或 $\ell \leq i + k - 1$（第一个
因子依归纳假设为 $0$），或 $\ell \geq i + k$，此时
$j \leq i + k \leq \ell$ 使第二个因子为零。当 $k = n$ 时
条件 $j \leq i + n - 1$ 对一切 $i, j \leq n$ 成立，故 $A^n =
0$。极小多项式整除 $X^n$，零化性质
确定了指数。
\end{proof}

\begin{theorem}[邓福德分解]\label{thm:b2:reduction:dunford}
设 $\chi_u$ 在 $K$ 上分裂（当 $K = \C$ 时自动成立）。则
存在\emph{唯一}的数对 $(d, \nu)$，满足
\[
u = d + \nu, \qquad d \text{ 可对角化}, \quad \nu
\text{ 幂零}, \quad d\nu = \nu d ,
\]
并且 $d$ 与 $\nu$ 都是
$u$ 的多项式。\index{邓福德分解}
\end{theorem}

\begin{proof}
\emph{存在性。}记 $\chi_u = \prod_{i=1}^{r} (X -
\lambda_i)^{m_i}$（$\lambda_i$ 互异），并令 $N_i = \ker(u -
\lambda_i)^{m_i}$ 称为\emph{特征子空间}。由
凯莱–哈密顿定理与核引理
（\cref{thm:b2:reduction:kernels}），
\[
E = N_1 \oplus \dots \oplus N_r ,
\]
其中投影 $\pi_i$ 为 $u$ 的多项式；每个 $N_i$ 都是
稳定的（$u$ 的多项式与 $u$ 可交换）。定义 $d = \sum_i \lambda_i
\pi_i$：$u$ 的多项式，可对角化（它以 $\lambda_i$
作用于 $N_i$ 上，故 $E$ 分解为其特征子空间）。于是 $\nu = u -
d$ 是 $u$ 的多项式（因而与 $d$ 可交换），且在每个
$N_i$ 上作用为 $u - \lambda_i$，且 $(u - \lambda_i)^{m_i} = 0$
在该处，$\nu^{\max m_i} = 0$ 于每个直和项上成立，故 $\nu$ 幂零。

\emph{唯一性。}设 $u = d' + \nu'$ 为另一个满足条件的数对。由于
$d'$ 与 $\nu'$ 彼此可交换，因而它们与 $u = d' +
\nu'$，从而与 $u$ 的每个多项式可交换 ——特别地与 $d$
及 $\nu$ 可交换。于是 $d - d'$ 可对角化（两个可交换的
可对角化映射可同时对角化：
见\cref{exo:b2:reduction:9}），且等于 $\nu' - \nu$，它是
幂零的：若 $\nu^k = 0$ 且 $\nu'^{k'} = 0$，可交换性
允许作如下的二项式展开
\[
(\nu' - \nu)^{k + k' - 1}
= \sum_{j=0}^{k+k'-1}\binom{k + k' - 1}{j}
\,\nu'^{\,j}\,(-\nu)^{k + k' - 1 - j} ,
\]
其中每一项皆为零：或 $j \geq k'$（第一个因子
为零），或 $k + k' - 1 - j \geq k$（第二个因子为零），且二者
之一恒成立。可对角化的幂零自同态为零（其
谱为 $\{0\}$，且在某组基下为对角形）：$d = d'$，
$\nu = \nu'$。
\end{proof}

\begin{example}[幂与指数]\label{ex:b2:reduction:powers}
$A = \begin{pmatrix} 3 & 1\\ -1 & 1\end{pmatrix}$：$\chi_A = X^2 -
4X + 4 = (X-2)^2$，只有单一特征值 $2$，特征子空间的
维数为 $1$，故不可对角化。邓福德分解：$D = 2I$，$N = A - 2I =
\begin{pmatrix} 1 & 1\\ -1 & -1\end{pmatrix}$，$N^2 = 0$。于是
\[
A^k = (2I + N)^k = 2^k I + k\,2^{k-1} N ,
\qquad
\eu^{tA} = \eu^{2t}(I + tN),
\]
分别由交换元的二项式定理，以及
在\cref{ch:b2:diffeq}中的指数级数对可交换的和项进行拆分得到。约化把
矩阵动力学化为标量动力学。
\end{example}

\begin{remark}[本卷中的几个视角]\label{rem:b2:reduction:perspectives}
约化是一个枢纽，下面列出值得关注的四个方面。在
\cref{ch:b2:nvs}中，适配范数把“所有特征值的模
$< 1$”转化为“某个算子范数 $< 1$”，使
谱决定幂与级数的收敛性。在
\cref{ch:b2:diffeq}中，把
\cref{ex:b2:reduction:powers}的做法化为
$X' = AX$ 的通解：邓福德分解将 $\eu^{tA}$ 拆分为
由多项式乘指数函数构成的块，稳定性则由
特征值的实部读出。在\cref{ch:b2:quadratic}中，
内积迫使得出纯线性代数所不能给出的结论：对称矩阵
成为\emph{正交}可对角化的，且完全没有
幂零部分。而在\cref{ch:b2:genfun}中，本章周末问题的
主导特征值渐近性态
以解析形式重新出现，表现为生成函数的最小奇点 ---
两种语言描述同一个增长率。
\end{remark}

\section{练习}

\begin{exercise}[$\star$]\label{exo:b2:reduction:1}
将下列矩阵对角化（求特征值、特征子空间的基以及可逆矩阵 $P$）：
\[
A = \begin{pmatrix} 1 & 2\\ 2 & 1 \end{pmatrix},
\qquad
B = \begin{pmatrix} 0 & 1 & 1\\ 1 & 0 & 1\\ 1 & 1 & 0
\end{pmatrix}.
\]
\end{exercise}

\begin{exercise}[$\star$]\label{exo:b2:reduction:2}
证明 $C = \begin{pmatrix} 1 & 1\\ 0 & 1\end{pmatrix}$ 不可对角化，
用特征子空间与用极小
多项式，两种方法各证一次。
\end{exercise}

\begin{exercise}[$\star$]\label{exo:b2:reduction:3}
设 $u$ 满足 $u^2 - 5u + 6\,\mathrm{id} = 0$。证明 $u$ 可对角化，
确定所有可能的谱，并把 $u^k$
写成 $\mathrm{id}$ 与 $u$ 的线性组合。
\end{exercise}

\begin{exercise}[$\star\star$]\label{exo:b2:reduction:4}
设 $u$ 可对角化，$F$ 为其不变子空间，证明
$u|_F$ 可对角化\emph{（对一个零化多项式作限制，
该多项式具有单的分裂根）}。
\end{exercise}

\begin{exercise}[$\star\star$]\label{exo:b2:reduction:5}
（斐波那契数列）设 $A = \begin{pmatrix} 1 & 1\\ 1 & 0\end{pmatrix}$。
将 $A$ 在 $\R$ 上对角化，并推出斐波那契数列的比内公式
（满足 $F_0 = 0$，$F_1 = 1$，$F_{n+1} = F_n +
F_{n-1}$）：
\[
F_n = \frac{\varphi^n - \psi^n}{\sqrt 5},
\qquad \varphi = \frac{1 + \sqrt5}{2},\ \psi = \frac{1 -
\sqrt5}{2}.
\]
\end{exercise}

\begin{exercise}[$\star\star$]\label{exo:b2:reduction:6}
设 $u \in \mathcal{L}(E)$，$u^2$ 可对角化，且 $u$
可逆（$K = \C$）。证明 $u$ 可对角化，
并在 $u$ 不可逆时给出一个反例。
\end{exercise}

\begin{exercise}[$\star\star$]\label{exo:b2:reduction:7}
对下列矩阵求邓福德分解、$A^k$ 及 $\eu^{tA}$：
\[
A = \begin{pmatrix} 2 & 1 & 0\\ 0 & 2 & 1\\ 0 & 0 & 2
\end{pmatrix}.
\]
\end{exercise}

\begin{exercise}[$\star\star$]\label{exo:b2:reduction:8}
设 $A \in \mathcal{M}_n(\C)$，满足 $A^k = I$，其中 $k \geq 1$。
证明 $A$ 可对角化，且其特征值为 $k$ 次
单位根。由此推出，任一有限阶可逆复矩阵若与
对角元全为1的三角矩阵相似，则必为单位矩阵。
\end{exercise}

\begin{exercise}[$\star\star\star$]\label{exo:b2:reduction:9}
（同时对角化）设 $u, v$ 为可对角化的
且可交换的。证明它们可\emph{同时}对角化，即
存在一组基使二者同时对角化。\emph{（$u$ 的每个特征子空间
均为 $v$ 不变；在其上将 $v$ 的限制对角化，
并利用\cref{exo:b2:reduction:4}。）}
\end{exercise}

\begin{exercise}[$\star\star\star$]\label{exo:b2:reduction:10}
设 $u \in \mathcal{L}(\C^n)$。证明 $u$ 可对角化，当且仅当
每个 $u$ 不变子空间都有一个 $u$ 不变的
补子空间。\emph{（对 $\Leftarrow$ 方向，取 $F = \sum_\lambda E_\lambda(u)$，
即全部特征子空间之和，并对其应用该性质；若存在非零的
不变补子空间 $G$，则对 $u|_G$ 作三角化便
会得到 $u$ 在 $G$ 中的一个特征向量，与 $G \cap
F = \{0\}$ 矛盾。）}
\end{exercise}

\begin{exercise}[$\star\star\star$]\label{exo:b2:reduction:11}
（谱半径，盖尔范德式简化版，$2\times2$ 情形的分析味道，
留待后文）设 $A \in \mathcal{M}_2(\C)$，其两个特征值的
模均 $< 1$。证明 $A^k \to 0$（当 $k \to \infty$ 时，逐元收敛）。
\emph{（三角化：$A = P(T)P^{-1}$，其中 $T$ 为上三角矩阵；
显式计算 $T^k$ --- 分别讨论特征值相等与
互异两种情形，并作估计。）}
\end{exercise}

\begin{exercise}[$\star\star$]\label{exo:b2:reduction:12}
设 $u \in \mathcal{L}(\C^n)$，满足 $\operatorname{rk} u = 1$（$n
\geq 2$）。证明 $\chi_u = X^{n-1}(X - \operatorname{tr} u)$，
并证明 $u$ 可对角化当且仅当 $\operatorname{tr} u
\neq 0$。\emph{（回顾 \cref{exo:b2:linalg:5}，有 $u^2 =
(\operatorname{tr} u)\,u$。）}
\end{exercise}

\section{问题：线性递推与友矩阵}

线性递推 $u_{n+k} = a_{k-1}u_{n+k-1} + \dots + a_0 u_n$
实为伪装的矩阵幂，约化可将其化为
闭式公式、增长率与误差估计。本周末
问题建立这一对应词典 --- 一侧为友矩阵，
另一侧为数列空间上的移位算子
--- 并证明\emph{线性递推的基本定理}
（一般解为 $\sum_i Q_i(n)\lambda_i^n$，求和遍及
特征多项式的根），并把由此得到的收获用于
$\sqrt2$ 的丢番图逼近、路径与单词的计数，
以及一组相互耦合的数列，它们只有
借助同时对角化才能解开。

\begin{problem}[{周末问题 --- 线性递推的
基本定理}]
\label{pb:b2:reduction:1}
取定 $k \geq 1$，标量 $a_0, \dots, a_{k-1} \in \C$，满足 $a_0
\neq 0$，首一多项式 $P = X^k - a_{k-1}X^{k-1} - \dots -
a_1 X - a_0$，以及递推式
\[
(\mathcal R)\colon\quad u_{n+k} = a_{k-1}u_{n+k-1} + \dots +
a_1 u_{n+1} + a_0 u_n \qquad (n \geq 0).
\]
$P$ 的\emph{友矩阵}为
\[
C =
\begin{pmatrix}
0 & 1 & & \\
 & \ddots & \ddots & \\
 & & 0 & 1\\
a_0 & a_1 & \cdots & a_{k-1}
\end{pmatrix}
\in \mathcal{M}_k(\C).
\]

\medskip
\noindent\textbf{第一部分 --- 友矩阵对应词典。}
\begin{enumerate}
  \item 证明数列 $(u_n)$ 满足 $(\mathcal R)$ 当且仅当
        向量 $v_n = (u_n, u_{n+1}, \dots,
        u_{n+k-1})^{\mathsf T}$ 满足 $v_{n+1} = Cv_n$，从而
        $v_n = C^n v_0$。
  \item 证明 $\chi_C = P$（沿第一列展开 $\det(XI - C)$，
        并对 $k$ 作归纳），再证 $\mu_C = P$
        亦成立\emph{（改用 $C^{\mathsf T}$，其中 $e_1$
        为循环向量，并注意矩阵与其转置有
        相同的极小多项式）}。
  \item 证明：设 $\lambda$ 为 $P$ 的任一根，则向量
        $(1, \lambda, \dots, \lambda^{k-1})^{\mathsf T}$ 张成
        $C$ 关于 $\lambda$ 的特征子空间；由此推出
        \emph{每一个} $C$ 的特征子空间维数均为 $1$，并
        证明 $C$ 可对角化当且仅当 $P$ 有 $k$
        个互异的根。
  \item 设 $P$ 有互异的根 $\lambda_1, \dots,
        \lambda_k$。证明几何数列
        $(\lambda_i^n)_n$ 构成解空间的一组基，
        故 $(\mathcal R)$ 的每个解都是 $u_n = \sum_i
        c_i\lambda_i^n$，其中常数 $c_i$ 唯一确定。
  \item 完全求解：$u_{n+2} = u_{n+1} + 6u_n$，$u_0 = 1$，
        $u_1 = 8$。
\end{enumerate}

\medskip
\noindent\textbf{第二部分 --- 移位算子与
基本定理。}设 $\mathcal{S}$ 为全体复数列构成的 $\C$ 上的
向量空间，并设 $S \in
\mathcal{L}(\mathcal{S})$ 为移位算子，$S\bigl((u_n)_n\bigr) =
(u_{n+1})_n$。
\begin{enumerate}[resume]
  \item 证明 $(\mathcal R)$ 的解集为 $\ker
        P(S)$，并证明其维数恰为 $k$ \emph{（把一个
        解映到其初始值）}。
  \item 说明为何核分解引理
        （\cref{thm:b2:reduction:kernels}）无需任何改动即可用于 $S$ 作用在
        无限维空间 $\mathcal{S}$ 上，
        并写出 $\ker P(S)$ 的相应分解，其中
        $P = \prod_{i=1}^{r}(X - \lambda_i)^{m_i}$（互异的
        $\lambda_i$ 均非零，因 $a_0 \neq 0$）。
  \item 对 $\lambda \neq 0$ 与 $m \geq 1$，证明
        \[
        \ker\,(S - \lambda\,\mathrm{id})^m
        = \bigl\{\,\bigl(Q(n)\,\lambda^n\bigr)_n : Q \in
        \C_{m-1}[X]\,\bigr\},
        \]
        且其维数为 $m$。\emph{（计算 $(S -
        \lambda)\bigl(Q(n)\lambda^n\bigr) =
        \lambda^{n+1}(\Delta Q)(n)$，其中 $\Delta Q = Q(X + 1) -
        Q(X)$，并利用 $\Delta$ 降低次数；
        求维数时借助初值证明其不超过 $m$。）}
  \item （线性递推的基本定理）由此得出：
        若 $P = \prod_{i=1}^{r}(X - \lambda_i)^{m_i}$，且
        $\lambda_i$ 互异且非零，则
        $(\mathcal R)$ 的解恰为如下形式的数列：
        \[
        u_n = \sum_{i=1}^{r} Q_i(n)\,\lambda_i^n,
        \qquad Q_i \in \C_{m_i - 1}[X],
        \]
        其中多项式 $Q_i$ 唯一确定。
  \item 完全求解：$u_{n+2} = 4u_{n+1} - 4u_n$，$u_0 =
        1$，$u_1 = 0$，并用 $u_2$ 检验答案。
\end{enumerate}

\medskip
\noindent\textbf{Part III --- 占优根与丢番图
红利。}
\begin{enumerate}[resume]
  \item 设各根均为单根，且 $\abs{\lambda_1} >
        \abs{\lambda_i}$（$i \geq 2$），$u_n = \sum_i c_i
        \lambda_i^n$，$c_1 \neq 0$。证明 $u_n \sim
        c_1\lambda_1^n$，且 $u_{n+1}/u_n \to \lambda_1$。
  \item （佩尔）定义 $a_{n+1} = a_n + 2b_n$，$b_{n+1} = a_n +
        b_n$，$a_0 = b_0 = 1$。证明 $q(a, b) = a^2 - 2b^2$
        满足 $q(a_{n+1}, b_{n+1}) = -q(a_n, b_n)$，从而
        $a_n^2 - 2b_n^2 = (-1)^{n+1}$；将此与
        矩阵 $M = \begin{psmallmatrix}1 & 2\\ 1 &
        1\end{psmallmatrix}$ 的行列式相联系。
  \item 由此推出误差估计
        \[
        \Bigl|\frac{a_n}{b_n} - \sqrt2\Bigr|
        = \frac{1}{b_n\,(a_n + \sqrt2\,b_n)}
        \leq \frac{1}{2b_n^2},
        \]
        并证明它以 $3 -
        2\sqrt2$ 为公比几何衰减，\emph{（求 $M$ 的特征值以及
        $b_n$ 的增长）}。
  \item （一般增长）由第 9 题证明：(a) 若每个
        根都满足 $\abs{\lambda_i} \leq \rho$，则
        $\abs{u_n} \leq C\,n^{m-1}\rho^n$，其中 $m = \max_i m_i$；
        (b) 若存在唯一的模最大根 $\lambda_1$，
        且 $Q_1 \neq 0$，则 $u_{n+1}/u_n \to
        \lambda_1$ --- 在第 10 题的解上检验此结论。
\end{enumerate}

\medskip
\noindent\textbf{Part IV --- 计数游走与词。}对于
一个有限图，其顶点集为 $\{1, \dots, N\}$，
\emph{邻接矩阵} $A$ 满足 $A_{ij} = 1$（若 $ij$ 是一条边），
否则为 $0$。
\begin{enumerate}[resume]
  \item 证明 $(A^n)_{ij}$ 是长为
        $n$ 的、从 $i$ 到 $j$ 的游走数（即由 $n$ 条边组成的序列，每一步沿一条边
        行走）。
  \item （三角形）对于 $3$ 个顶点的完全图，
        $A = J - I$：利用 $J$ 的谱
        （见\cref{ex:b2:linalg:onesmatrix}），证明
        \[
        (A^n)_{ii} = \frac{2^n + 2(-1)^n}{3},
        \qquad
        (A^n)_{ij} = \frac{2^n - (-1)^n}{3} \quad (i \neq j),
        \]
        并通过列举游走，在 $n = 2$ 时检验两式。
  \item （不含 $11$ 的词）设 $w_n$ 为
        长为 $n$ 且不含两个连续 $1$ 的二进制词的个数，
        用末位字母对词编码，得到转移矩阵，
        证明 $w_{n+2} = w_{n+1} + w_n$，由此推出 $w_n =
        F_{n+2}$（斐波那契数，见\cref{exo:b2:reduction:5}），并
        求增长率 $\lim w_{n+1}/w_n$。
  \item （路径）对于路径图 $1 - 2 - 3$，证明
        $A$ 的特征值为 $\sqrt2, 0, -\sqrt2$，
        对应特征向量为 $(1, \pm\sqrt2, 1)$ 与 $(1, 0, -1)$，并由此推出
        长为 $n$ 的、从一端到另一端的游走数为
        $\bigl((\sqrt2)^n + (-\sqrt2)^n\bigr)/4$：
        当 $n$ 为奇数时为零，$2^{\,n/2 - 1}$ 适用于偶数 $n$。
        在 $n = 4$ 时检验。
  \item （迹公式）证明长为
        $n$（起点任意）的闭游走总数为
        $\operatorname{tr}(A^n) = \sum_i \lambda_i^n$，并
        在三角形上验证此式。
\end{enumerate}

\medskip
\noindent\textbf{Part V --- 数列的环：同时
对角化。}固定 $k \geq 3$，设 $\omega = \eu^{2\iu\pi/k}$，
        并设 $W \in \mathcal{M}_k(\C)$ 为循环位移：$W e_i =
e_{i+1}$（下标按模 $k$，列编号依次为 $0, \dots, k-1$）。
\begin{enumerate}[resume]
  \item 证明 $W^{\mathsf T}$ 为
        $X^k - 1$ 的友矩阵，由此推出 $\chi_W = \mu_W = X^k - 1$，且 $W$
        可对角化，其 $k$ 个单特征值为
        $\omega^j$，特征向量为 $f_j = (1, \omega^{-j},
        \omega^{-2j}, \dots, \omega^{-(k-1)j})^{\mathsf T}$。
  \item \emph{循环矩阵}是形如 $C = c_0 I + c_1 W + \dots
        + c_{k-1}W^{k-1}$ 的矩阵。证明所有循环矩阵两两可交换，
        基 $(f_0, \dots, f_{k-1})$ 同时对角化
        \emph{全部}循环矩阵，并证明
        $C$ 的特征值为 $\widehat c(\omega^j) = \sum_m
        c_m \omega^{jm}$，$j = 0, \dots, k-1$。
  \item 由此推出 $\det C = \prod_{j=0}^{k-1} \widehat
        c(\omega^j)$，并验证 $k = 3$ 时可得
        \cref{exo:b2:linalg:8} 中的分解式。
  \item （项链平均）设 $x^{(n+1)} = Mx^{(n)}$，其中
        $M = \frac12(W + W^{-1})$：排成环状的 $k$ 个数中，每个数
        都被替换为其两个相邻数的平均值。证明
        $M$ 的特征值为
        $\cos(2\pi j/k)$，并证明 $x^{(0)}$ 在
        $f_0$ 上的系数即为均值 $\frac1k\sum_m x^{(0)}_m$
        \emph{（对 $f_j$ 的各坐标求和）}。
  \item 总结：当 $k$ 为奇数时，$x^{(n)}$ 收敛于
        常向量，其值为初始值的均值；
        当 $k = 4$ 时，指出导致不收敛的特征值，
        以及精确的障碍
        （一个必须为零的交错均值系数）。
  \item （综合）各用一句话说明：友矩阵如何
        将 $(\mathcal R)$ 的分析化为
        约化；核分解引理在何处
        不需要有限维；占优特征值为何
        支配增长率与丢番图误差；为何
        邻接矩阵的幂可以计数游走；以及
        可交换矩阵有何好处。指出两座高峰：线性递推的基本
        定理，以及 --- 针对第 IV 部分的正矩阵，
        见本科三年级分册 --- 即
        佩龙--弗罗贝尼乌斯定理。
\end{enumerate}
\end{problem}
